\section{Introduction}
\label{sec-introduction}

Calibrer un modèle, c'est choisir ses paramètres pour que les prix qu'il donne collent aux prix du marché. Pour le
modèle de \citet{heston1993}, il n'existe pas de formule fermée des prix : chaque surface de volatilité implicite
se calcule par une méthode de Fourier, qui prend quelques millisecondes. Un optimiseur demande des centaines
d'évaluations de la surface, et une banque recalibre ses modèles plusieurs fois par jour, sur de nombreux sous-jacents :
le pricer devient le goulot d'étranglement.

\citet{horvath2021} proposent de séparer le travail en deux temps. Hors ligne, un réseau de neurones apprend la
fonction qui associe aux paramètres du modèle la surface de volatilité sur une grille fixe de maturités et de
strikes ; les exemples d'entraînement sont produits une fois pour toutes avec le pricer exact. En ligne, la
calibration utilise le réseau à la place du pricer : une évaluation n'est plus qu'un passage avant dans un petit
réseau, beaucoup plus rapide. Le coût est déplacé vers la génération des données et l'entraînement, qui ne sont
faits qu'une fois.

Nous appliquons cette méthode au modèle de Heston et cherchons à répondre à quatre questions :
\begin{itemize}
  \item quelle architecture choisir, sachant que le réseau doit être à la fois précis et au moins dix fois plus
  rapide que le pricer Fourier ;
  \item comment l'entraîner (optimiseur, initialisation, fonction d'activation) ;
  \item un réseau de neurones fait-il mieux que des méthodes d'apprentissage classiques sur ce problème ;
  \item la calibration par le réseau est-elle aussi bonne que la calibration par Fourier, sur des surfaces
  simulées et sur le marché (SPX), et à partir de combien de calibrations le temps d'entraînement est-il
  rentabilisé.
\end{itemize}

La \cref{sec-donnees} présente le modèle, le pricer de référence, le jeu d'entraînement et les surfaces de marché.
La \cref{sec-mise-en-oeuvre} décrit l'organisation du code. Les \cref{sec-architectures,sec-entrainement}
choisissent l'architecture puis l'entraînement du réseau. La \cref{sec-dl-vs-ml} le compare à des méthodes
classiques, et la \cref{sec-calibration} l'utilise pour calibrer. La \cref{sec-conclusion} conclut.
